\documentclass[a4paper,12pt]{article}
% Paquetes de la ams
\usepackage{amsmath,amsthm,amssymb,amsfonts}
% Posibilidad de mover la pagina
\usepackage[a4paper]{geometry}
% Saco la indentacion en todos los parrafos.
\usepackage{parskip}
% Codificacion UTF-8
\usepackage[utf8]{inputenc}
% Tablas e imagenes en espaniol
\usepackage[spanish,es-tabla]{babel}
% Mejores graficos
\usepackage{graphicx}
% tablas mas lindas
\usepackage{booktabs}
% Posibilidad de tocar los encabezados
%\usepackage{fancyhdr}
%\pagestyle{fancy}
% Posibilidad de meter subfigurasw
%\usepackage[font=footnotesize, labelfont=it]{subcaption}
% Links a urls
\usepackage{url}
% Linkear referencias en pdfs
\usepackage{hyperref}
% Texto mas lindo para los pie de figura
\usepackage[margin=10pt,font=small,labelfont=bf, labelsep=endash]{caption}
% Mejores autores
\usepackage[affil-it]{authblk}
% Compatibilidad con PDF/A
\usepackage{xmpincl}
% Hoja a4 mas ancha
\usepackage{a4wide}
% Citas
%\usepackage[backend=biber,style=ieee]{biblatex}
%\addbibresource{biblio.bib}
% Cambio and por y
\renewcommand\Authand{y }
\renewcommand\Authands{, y }

% Codigo
\usepackage{listings}

% Coloreo los links
\usepackage[usenames,dvipsnames]{xcolor}
\hypersetup{colorlinks,
     linkcolor={red!50!black},
     citecolor={blue!50!black},
     urlcolor={blue!80!black} }
% Graficos con tikz
\usepackage{tikz}

% Dir tree
\usepackage{dirtree}

% Pagina en blanco cuando ha
\usepackage{emptypage}

% Subfiguras
\usepackage{subfig}

% Tablas largas
\usepackage{longtable}
% Fechas automáticas
\usepackage{advdate}    % Advancing/saving dates
\usepackage{datetime}   % Dates formatting
\usepackage{datenumber} % Counters for dates

\newdateformat{mydate}{\THEDAY\ de \shortmonthname[\THEMONTH]}
\renewcommand\AdvanceDate[1][\@ne]{\global\advance\day#1 \FixDate}

% Menues y path
\usepackage[os=win]{menukeys}
\renewmenumacro{\directory}{pathswithfolder}
\definecolor{A11}{HTML}{B2DF8A}
\definecolor{A12}{HTML}{33A02C}
\definecolor{A23}{HTML}{FDBF6F}
\definecolor{A24}{HTML}{FF7F00}
\definecolor{B15}{HTML}{FB9A99}
\definecolor{B16}{HTML}{E31A1C}
\definecolor{B27}{HTML}{A6CEE3}
\definecolor{B28}{HTML}{1F78B4}

% Ejemplos, observaciones y teorema
\theoremstyle{definition}
\newtheorem{exa}{Ejemplo}[section]
\newtheorem*{obs}{Observación}
\newtheorem{que}{}[section]
\newtheorem{dex}{Definicion}[section]

\usepackage{float}
\usepackage{enumitem}
%\includeonly{guia1}
\graphicspath{{./figs/}}
% Renombro lo capitulos
\addto\captionsspanish{\renewcommand{\chaptername}{Clase}}


\title{ {\large Nivel 2} \\ Introducción al color del Oceano}
\author{}

%\author[*]{Laura Rouco}


%\titlepic{\vfill\includegraphics[height=1.5cm]{../figs/logo-ministerio.png}\hfill\includegraphics[height=1.5cm]{../figs/logo-conae.png}}
\pagenumbering{gobble}
\date{}

\begin{document}
%\frontmatter
\maketitle
%\titlepage
%\tableofcontents
%\chapter{Prefacio}

INTRODUCCION A LA TECNOLOGIA SAR. OBJETIVOS DEL CURSO. ETC.

\vskip-2.2cm
%\mainmatter
El Curso Nivel 2: Introducción al color del oceano en su Modalidad Virtual está destinado a usuarios de imágenes satelitales que hayan realizado y aprobado el Curso Nivel 1: Introducción a Teledetección o tener conocimientos equivalentes al mencionado curso en teledetección. El objetivo de esta capacitación es que los participantes adquieran los principales conceptos y aplicaciones de las imágenes de color del oceano. 
\section*{Organización general}
\begin{itemize}
    \item Modalidad: a distancia, a través de la plataforma virtual de la UEFM.
    \item Duración: 5 clases.
    \item Carga horaria: 25 horas.
    \item Requisitos: Haber realizado y aprobado el Curso Nivel I: Introducción a Teledetección o tener conocimientos equivalentes al mencionado curso en teledetección. Contar con una PC con Windows 7 o superior, Linux. Procesador i3 o superior, 8GB de memoria RAM, 30GB de espacio en disco disponible y una conexión a internet de 1 Mbit/s
        
\end{itemize}

\section*{Contenidos}
\begin{enumerate}[leftmargin=*, label={Clase \arabic* -}]
    \item \emph{Introducción:}
    Sensores remotos en color del oceano. Satélites y órbitas. Satélites geosincronicos. Satélites heliosincronicos. Técnicas de adquisición de datos. Niveles de procesamiento. Archivos de datos de color del oceano.
    \item \emph{Radiación electromagnética:}
    Introducción. Ondas electromagneticas. Espectro electromagnético. Descripción de radiación. Magnitudes físicas.
    \item \emph{Propiedades de la atmósfera:}
    Descripción de la atmósfera. Moleculas: emisión y absorción. Scattering. Atenuación en la atmósfera. Ecuación de transferencia radiativa. Soluciones a la ecuación de transferencia radiativa. Transmitacia difusa.
    \item \emph{Color del océano:}
    Absorción y scattering del phitoplackton, particulas y material disuelto. Instrumentos usados en color del oceano. Correcciones atmosféricas. Calibración y validación. Modelos empíricos y semiempíricos. 
    \item \emph{Misión SABIA-Mar y aplicaciones:}
    Plan Espacial Nacional. Misión SABIA-Mar. Misiones de color del oceano. Aplicaciones de los datos de color del oceano.
\end{enumerate}


\end{document}
